\begin{helper}
Let $\Delta>0$, let $0<\lambda\le 1$, and let $K$ be the finite pair
kernel contribution for a fixed independent pair.  Fix a split parameter
$\alpha\le 1$.  Suppose the compact-window
estimates from the finite pair integral give the following two lower bounds:
if $\alpha\le\lambda$, then
\[
K\ge
\frac{1}{\Delta^2}\left(\frac{2}{1+\lambda}-6\alpha\right),
\]
while if $\lambda<\alpha$, then
\[
K\ge
\frac{1}{\Delta^2}\left(\frac53+2\alpha\right).
\]
Then, in both cases,
\[
K\ge
\frac{1}{\Delta^2}
+\frac{1}{3\Delta^2}\lambda
\left(\frac{2}{\lambda(1+\lambda)}-1\right)
-\frac{6\alpha}{\Delta^2}.
\]
Equivalently, the finite pair contribution is at least the correction term
\[
\frac{1}{\Delta^2}
+\frac{1}{3\Delta^2}\lambda
\left(\frac{2}{\lambda(1+\lambda)}-1\right)
\]
up to the explicit loss $6\alpha/\Delta^2$.
\end{helper}
\begin{proof}
We argue by the two cases in the compact-window split.

First suppose $\alpha\le\lambda$.  The low-overlap compact-window estimate
gives
\[
K\ge
\frac{1}{\Delta^2}\left(\frac{2}{1+\lambda}-6\alpha\right).
\]
It is therefore enough to prove
\[
1+\frac13\lambda\left(\frac{2}{\lambda(1+\lambda)}-1\right)
\le \frac{2}{1+\lambda}.
\]
Since $0<\lambda$, the left side simplifies to
\[
1+\frac13\left(\frac{2}{1+\lambda}-\lambda\right).
\]
Subtracting it from $2/(1+\lambda)$ gives
\[
\frac{(1-\lambda)^2}{3(1+\lambda)}\ge0,
\]
because $1+\lambda>0$.  Dividing by $\Delta^2>0$ and subtracting
$6\alpha/\Delta^2$ proves the desired lower bound in this case.

Now suppose $\lambda<\alpha$.  The high-overlap fixed-window estimate gives
\[
K\ge
\frac{1}{\Delta^2}\left(\frac53+2\alpha\right).
\]
For every $0<\lambda\le1$,
\[
1+\frac13\lambda\left(\frac{2}{\lambda(1+\lambda)}-1\right)
\le \frac53.
\]
Indeed, after multiplying by the positive number $3(1+\lambda)$, this is
equivalent to
\[
3(1+\lambda)+2-\lambda(1+\lambda)\le 5(1+\lambda),
\]
which reduces to $\lambda(1-\lambda)\ge0$.  Hence
\[
\frac{1}{\Delta^2}
+\frac{1}{3\Delta^2}\lambda
\left(\frac{2}{\lambda(1+\lambda)}-1\right)
-\frac{6\alpha}{\Delta^2}
\le
\frac{5/3-6\alpha}{\Delta^2}.
\]
Since $\lambda<\alpha$ and $0<\lambda$, we have $\alpha>0$, so
$5/3-6\alpha\le 5/3+2\alpha$.  Dividing by the positive number
$\Delta^2$ and using the high-overlap estimate completes the proof.
\end{proof}
